\documentclass[../main.tex]{subfiles}
\begin{document}
%==========================================TIÊU ĐỀ TẬP========================================
\begin{table}[h]
  \begin{adjustwidth}{-1cm}{}
    \begin{tabular}{>{\centering\arraybackslash}p{6cm}|>{\raggedright\arraybackslash}p{12.5cm}}
      \multirow{5}{6cm}
      % Cột bên trái: ÉPISODE và số tập
      {\\[-40pt]
        \flaregothic\fontsize{20pt}{20pt}\selectfont \centering
        \color{eptype}ÉPISODE\color{black}\\
        \burbank\fontsize{72pt}{72pt}\selectfont
        \color{epnum}133\color{black}\\[6pt]
        \cafeteria\fontsize{20pt}{20pt}\selectfont\color{epnum}(10-12)\color{black}
        \\[18pt]
        % \flaregothic\fontsize{14pt}{14pt}\selectfont
        % \color{epattr}BIENVENUE, \uline{\color{Banpire} Mel} !
        % \flaregothic\fontsize{18pt}{18pt}\selectfont
        % \color{epattr}ÉPISODE PERDU
        \color{black}
      }
       &
      % Dòng thứ nhất: tên tập tiếng Nhật
      {\fotskip\fontsize{18pt}{18pt}\selectfont \ruby{友}{とも}\ruby{達}{だち}のこと\ruby{騙}{だま}してきたわ
      }          \\[6pt]
      % Dòng thứ hai: tên tập tiếng Anh
       & {\cafeteria\fontsize{18pt}{18pt}\selectfont Hoodwinked, Bamboozled, Led Astray}             \\[4pt]
      % Dòng thứ ba: tên tập tiếng Việt
       & {\vnmsans\fontsize{18pt}{18pt}\selectfont Miếng dán ma thuật và những con rối nhân bản}     \\[8pt]
      % Dòng thứ tư: Nhân vật xuất hiện
       & {\sen{\fontsize{14pt}{14pt}\selectfont Track used: Alone Remix – Alan Walker}}
      \\[6pt]
      % Dòng cuối cùng: Thời gian diễn ra của tập (thường sẽ là ngày phát hành)
       & {\continuum\fontsize{16pt}{16pt}\selectfont Ngày phát hành: 28/11/2021}                     \\[8pt]
       & (dựa theo bản dịch của \href{https://www.facebook.com/mamisubteam}{MamiSubs - Mami Fansub}) \\[18pt]
    \end{tabular}
  \end{adjustwidth}
\end{table}
%=========================================TRUYỆN CHÍNH========================================
\color{black}

Một buổi tối cuối tháng Mười Một. Những ngày cuối thu se lạnh, khi ánh nắng chiều nhạt dần và gió bắt đầu mang theo hơi thở của mùa đông đang tới gần. Ở Tokyo ảo, đường phố đã thấp thoáng những cửa hiệu treo đèn Giáng sinh từ sớm, còn trong các văn phòng lớn, một số người đang tan ca, số khác thì lặng lẽ ở lại làm việc thêm chút nữa.

Trong một buổi phát trực tiếp của Shion, nàng Phù thủy hớn hở khoe với người xem một phát minh mới mà cô đã âm thầm chế tạo trong thời gian gần đây: ``Mọi người nghe này\~{}! Mình vừa hoàn thành một công cụ ma thuật mới, gọi là miếng dán thần kỳ! Khi dán cái này lên ai đó, mình sẽ điều khiển được cơ thể của họ\~{}! Dễ thương đúng không?''

\img{10-12a}

Ngồi trong phòng riêng, cậu Tổng chăm chú theo dõi buổi phát ấy trên màn hình laptop, một tay gác lên trán.

``\dots Chưa biết nó có hiệu quả thật không\dots\ mà thực ra, mình còn chẳng hiểu rõ nó là cái quái gì,'' Kuon thầm nghĩ, mắt vẫn không rời khỏi hình ảnh Shion đang lăng xăng trên màn hình.

Cùng lúc đó, ở khu làm việc bên cạnh, Korone cũng đang xem buổi stream ấy, đôi mắt sáng lấp lánh nhưng gương mặt lại không biểu lộ một chút cảm xúc nào. Cô ngồi bệt trên sàn, cằm chống lên gối, chẳng ai đoán được trong đầu cô đang tính toán điều gì.

\img{10-12b}

Có thể là một trò đùa mới?

Một món ``đồ chơi'' để thử trên Okayu chăng?

Hay là một thứ gì đó hay ho hơn nữa?

``Ồ\~{}?'' Korone bật ra một tiếng nhẹ, ánh mắt hứng thú. Đôi tai chó khẽ giật giật.

Cô không biết rằng\dots\ chính mình sẽ là một phần trong trò chơi kỳ lạ sắp tới.

\begin{center}
  \uline{Sáng hôm sau.}
\end{center}

Văn phòng chính tầng ba còn khá vắng. Kuon và Shion là hai người đến sớm nhất, như thường lệ. Dường như biết rằng sẽ chưa ai tin được miếng dán đó thực sự công hiệu, cô đã rủ cậu Tổng tới để biểu diễn và đồng thời làm chứng cho người xem rằng cô không nói dối.

Cầm miếng dán ma thuật trên tay, cậu không khỏi nhíu mày nghi ngờ: ``Ờm\dots\ Em có chắc là cái này hoạt động thật không đấy?''

``Em biết là anh sẽ không tin mà,'' nàng Phù thủy vênh mặt, ``vậy để em cho anh thấy\dots\ Shion phiên bản 2.0 đây!''

Cô nàng hớn hở trưng ra một chiếc ma-nơ-canh hình người, hay đúng hơn là một con rối ma thuật do chính cô tạo ra từ trước. Nó có hình dạng giống hệt cô, từ mái tóc xám tím cho đến bộ váy quen thuộc.

``Mà sao nó đứng trơ ra thế kia? Không phản ứng gì à?''

``Dĩ nhiên rồi. Nó chỉ là một cái xác thôi, cho đến khi\dots\ này nhé!''

Nói rồi, cô mở ngăn kéo, lấy ra một chiếc hộp nhỏ, trong đó có nhiều miếng dán hình tròn, vẽ đầy ký hiệu ma thuật. Cô bóc một miếng, dán vào sau gáy con rối, rồi dùng một cây bút phép vẽ một đường xoắn lạ mắt từ tâm miếng dán lan ra các khớp vai và gáy.

Ngay lập tức, con rối giật mạnh một cái, rồi bắt đầu cựa quậy. Chỉ trong vài giây, nó đã chuyển sang trạng thái hoạt động linh hoạt, có phần hơi quá mức.

\begin{figure}[H]
  \centering
  \includegraphics[width=12.5cm]{../../../../Image/Book?/10-12c.png}
  \caption{Shion hàng ``pha-ke'' (bên trái) vươn vai lên cao, cử động như một con người thực sự.}
\end{figure}

Nó đứng thẳng dậy, khom lưng, vặn mình, rồi bắt đầu lặp lại các động tác trong bài nhảy mà Shion thường biểu diễn trên sân khấu. Kỳ lạ hơn nữa, nó còn thực hiện linh hỏa hơn cả bản gốc, khiến cho cả chủ nhân của con rối cũng bất ngờ.

Nàng Phù thủy cười đắc chí: ``Thấy chưa nào\~{}? Giờ anh đã tin chưa?''

Cậu Tổng đưa tay chỉnh lại chiếc mình đang há hốc lại, rồi lắp bắp trả lời, dù rằng cậu vẫn không thể tin được: ``C-Cái đó\dots\ Ừ. Anh tin. Đúng là có tác dụng thật.''

``He he\~{} Em đúng là một thiên tài!'' cô nàng chống hông, cười như thể cô đã thực hiện thành công thí nghiệm sau nhiều năm tìm hiểu.

Kuon rõ ràng không thoải mái với cái tính tự đại của cô, liền lên tiếng nhắc nhở: ``Vẫn cái kiểu huênh hoang chưa thèm sửa kìa.''

Nàng Phù thủy thoáng tụt hứng đôi chút, phồng má lên bực bội: ``Ơ kìa! Thành công rồi thì phải khoe chứ!''

Con rối tiếp tục thể hiện những động tác cực khó: nâng chân lên ngang mặt, cúi người xuống kiểu Jack-O Challenge, rồi còn múa tay theo kiểu hoạt hình.

``Công nhận miếng dán này đa năng thật đấy.''

``Ờ, đa năng,'' cậu Tổng thở dài một hơi, rồi buông ra một yếu điểm trước mắt của nó, ``mỗi tội không biết nói.''

Shion khoanh tay lại, gật gù đồng tình với cậu. ``Cũng phải\dots\ Đây vẫn là bản thử nghiệm thôi. Nhưng nếu ta truyền âm thanh vào trong\dots''

Cô rút từ túi ra một chiếc loa mini, chuẩn bị lắp vào bên trong con rối thì\dots

*RẦM!*

Cánh cửa văn phòng bất ngờ bật mở. ``\textbf{Yubi yubi\~{}!!}''

Một giọng nói lanh lảnh vang lên từ hành lang. Cả hai giật nảy, quay lại. Không cần nhìn kỹ, họ cũng biết ngay ai vừa bước vào.

Shion thì lại bỗng sáng mắt như vừa nảy ra một ý tưởng táo bạo. ``À, đúng rồi!''

``Gì vậy?'' Kuon rướn mày, rõ ràng là đoán được ý định của đối phương. ``Đừng nói với anh là em-''

``Suỵt. Im nào. Đi theo em, rồi anh sẽ thấy!''

``Khoan, gì cơ? Ê này, đợi đã!!''

Không kịp phản ứng, cậu bị cô nàng kéo đi như một cục đồ chơi, còn bản thân nàng Phù thủy đã bắt đầu chuẩn bị cho màn thử nghiệm mới.

Bởi vì bản sao ki không phải là bản sao duy nhất.

Và người vừa mới vào phòng, chính là mục tiêu tiếp theo trong kế hoạch nhân bản kỳ quái này.

\ct

Korone sán lại, giơ tay vẫy chào đầy năng lượng: ``Chào buổi sáng, Shion-chan\~{}!''

Bản sao của nàng Phù thủy (hiện tại vẫn đang được điều khiển bằng phép thuật) cũng quay người lại, cười y hệt bản gốc: ``Chào buổi sáng, Korone\~{}!''

\begin{figure}[H]
  \centering
  \begin{subfigure}[t]{0.45\textwidth}
    \centering
    \includegraphics[width=8cm]{../../../../Image/Book?/10-12d1.png}
  \end{subfigure}
  \quad
  \begin{subfigure}[t]{0.45\textwidth}
    \centering
    \includegraphics[width=8cm]{../../../../Image/Book?/10-12d2.png}
  \end{subfigure}
\end{figure}

Nàng Khuyển hơi nheo mắt, tay chống hông nhưng vẫn không giấu khỏi niềm thích thú: ``Nè~ Đừng có bắt chước tôi chứ! Thô lỗ lắm đó!''

``Đâu có\~{} Tôi chỉ chào lại thôi mà\~{}'' bản sao của nàng Phù thủy đáp lại tỉnh bơ.

Trong khi đó, trên trần nhà văn phòng, hở ra một cái lỗ trần nhà, để lộ đôi mắt đang lấp lánh của nàng ``Tỏi'' và khuôn mặt đổ mồ hôi của cậu Tổng. Cả hai nấp trong không gian chật hẹp, theo dõi tình hình bên dưới qua một khe hẹp.

\img{10-12e}

Shion nhỏ giọng, nhe răng cười gian: ``\hl{Heh. Cậu ta không nhận ra gì luôn. Korone đúng là khờ thật\~{}}''

``\hl{Chơi bài kiểu\dots{} chọc chó à?}'' Kuon nhíu mày, cặp mắt vẫn dán chặt vào Korone đang cười tít mắt. ``\hl{Anh không nghĩ em ấy dễ bị-}''

Cô nàng lập tức cắt ngang lời cậu: ``\hl{Khẽ nào! Cậu ta đang chăm chú lắm đấy-}''

Cô chưa kịp dứt câu, cả Shion và Kuon hoảng hốt nghiêng người nấp sang bên, nín thở như bị đóng băng. Ngay lúc này, nàng Khuyển ngước mắt lên, chẳng nói lời nào, cặp mắt nâu của cô ánh lên ánh nhìn vô hồn kỳ lạ, kiểu ``mắt cá chết'' đặc trưng khiến ai nhìn cũng phải lạnh gáy.

\img{10-12f}

Phải vài giây sau khi Korone quay đầu trở lại và tiếp tục trò chuyện với bản sao, cả hai mới dám thở ra.

Cậu Tổng lau mồ hôi trên trán, thở mạnh: ``\hl{Vãi\dots{} Nguy hiểm thật\dots}''

Nàng Phù thủy che miệng cậu, suỵt một tiếng, rồi cười một nham hiểm như một đứa trẻ tìm được niềm vui riêng: ``\hl{Mà\dots{} công nhận, trò này vui phết\dots}''

Kuon bật cười một tiếng rất bé, thở dài một tiếng vì không biết nên hùa theo hay nên quở trách của cô. Cậu ngước mắt xuống, thầm nhớ lại lúc cậu nhìn lũ trẻ ở quê: ``\hl{Anh có nhớ lúc trước còn nhìn đám trẻ chọc chó rồi chạy tán loạn nữa\dots}'' nhưng rồi cũng nghiêm giọng cảnh cáo: ``\hl{Cẩn thận không chó cắn đấy nha.}''

``\hl{Em biết cách trị cậu ta rồi\~{} Cứ để đó cho em!}'' cô giơ ngón tay lên, đọc nhẩm câu thần chú rồi tiếp tục điều khiển con rối của mình.

Bên dưới, bản sao của Shion bắt đầu nhảy múa loạn xạ: từ điệu can-can, đến hip-hop, rồi chuyển qua vài động tác rô-bốt. Korone, vẫn chưa mảy may nghi ngờ, vỗ tay không ngớt: ``Uoaa\~{} Ngầu quá Shion-chan! Bà luyện kiểu gì mà dẻo vậy luôn á!''

Shion và Kuon trên trần nhìn nhau, rồi cùng lắc đầu, với một người cảm thấy tội, một người lại không giấu nổi vẻ thích thú.

``\hl{Koro-chan ngờ nghệch đến thế à?}'' cậu Tổng lẩm bẩm, ``\hl{Tưởng em ấy khôn lắm chứ\dots}''

``\hl{Công nhận luôn,}'' nàng Phù thủy đáp, thở ra ngáp một hơi. ``\hl{Mà nếu thế này mãi cũng chán\dots}''

Cả hai vừa nói xong thì\dots\ Korone lại ngước mắt lên khe trần lần nữa. Cái nhìn chằm chằm vô cảm của cô khiến họ rùng mình, như thể Korone đã biết rõ từ đầu.

\img{10-12f2}

Nàng Khuyển ngước nhìn bản sao của Shion lần nữa, tới lúc này họ mới thở phào nhẹ nhõm.

``\hl{\dots Hạ màn thôi nhỉ?}'' cô nàng hỏi.

``\hl{Ừ. Đi đêm lắm có ngày gặp ma chứ chả đùa,}'' cậu Tổng đáp lại, dần lùi ra lỗ thông gió. ``\hl{Anh sẽ vào bằng cửa trước như mọi ngày.}''

``\hl{Vậy anh cứ cư xử bình thường nhé.}''

Kuon đáp lại cô bằng một cái gật đầu nhẹ, rồi bắt đầu leo từ thang xuống.

Dưới kia, Korone vẫn mải mê quan sát bản sao của bạn mình, mắt sáng rực như đèn pin. ``Tiếp theo bà định làm gì nữa đây\~{}?''

Ngay lúc ấy, từ trần nhà, một bóng người nhảy xuống. Cùng với cú tiếp đất là tiếng hô vang: ``\textbf{Tada\~{}!! Bà đã bị lừa rồi nha!!}''

\img{10-12g}

Tay cô cầm tấm bảng đỏ chót, chữ to đùng đập vào mắt Korone. Trong khi nàng chó vẫn chưa kịp phản ứng, cửa phòng lại bật mở.

``Chào buổi sáng mấy đứa- Ô vãi! Gì đây!?''

Cả căn phòng giờ có hai Shion đứng song song. Korone giật mình nhìn qua nhìn lại. ``Hở? C-Cái gì thế này? Ý bà là sao đây?''

Dẫu rằng Kuon biết mọi thứ từ trước, cậu vẫn giả vờ ngạc nhiên để không bị cô nghi ngờ: ``S-Sao lại có\dots\ hai Shion-chan thế kia''

Nàng Phù thủy (thật) quẳng tấm bảng đi rồi đáp lại bằng một nụ cười vênh váo: ``Từ nãy tới già bà đang nói chuyện với con rối của tôi á\~{}''

``Ể!? C-Cái gì cơ!?'' / ``Thật á?'' cả nàng Khuyển và cậu Tổng đều bất ngờ với lời giải thích của cô.

Shion tiến lại gần, đặt nhẹ tay lên bản sao của mình. Lập tức, con rối đứng khựng lại, ngừng hoạt động hoàn toàn, như một bức tượng sáp.

``Y như một con rô-bốt kìa\dots'' cậu Tổng há hốc miệng, choáng ngợp với công cụ lợi hại này.

Korone thì vẫn vỗ tay không ngớt như một đứa trẻ vừa được xem ảo thuật: ``Oaaa\~{} Hay ghê á!''

Nhận được lời khen từ cả hai người, nàng Phù thủy nhỏ giương khuôn mặt tự mãn, hai tay chống hông, bừng bừng khí chất của một nhà phù thủy tài ba: ``Hưm hưm\~{} Tuyệt không nào? Tài năng của tôi mà lị!''

``Gì mà tự tin thế\dots'' nàng Khuyển vỗ vai, cho rằng cô ấy hơi tự đại.

Cậu Tổng khoanh tay, đáp lại: ``Anh thấy em đang nổ quá đấy. Tém tém lại đi.''

Shion lượn ra cửa, hí hửng như vừa trúng số. ``Vui thật đấy! Nhớ theo dõi buổi stream tối nay nha, em sẽ kể vụ này đấy!''

Tưởng mọi chuyện kết thúc ở đó. Nhưng khi cô quay đi, ánh mắt Kuon lặng lẽ bám theo, chợt thấy một điều bất thường.

``K-Khoan đã\dots\ Có gì đó không đúng\dots''

Ở sau gáy của cô nàng lấp ló một miếng dán y hệt miếng ban nãy. Trước khi cậu kịp lên tiếng, Korone đã giơ một ngón tay, nhẹ nhàng quẹt không khí như vẽ lệnh, miệng lẩm bẩm: ``Vẫn chưa đến lúc hạ màn đâu\~{}'' Nàng Phù thủy đứng khựng lại, giơ hai tay lên cao như vừa bị trúng bùa.

\img{10-12h}

Cậu con trai sững người, mồ hôi đổ ra đầm đìa: ``Em ấy phát hiện ra rồi ư!? Mình biết mà\dots\ không tài nào qua mắt được em ấy đâu\dots''

Trong lúc cậu vẫn còn đang chìm trong suy nghĩ hỗn loạn vì cuối cùng bị đối thủ bắt bài, một tiếng gọi trong phòng đưa cậu về thực tại. ``\dots-senpai? Kuon-senpai ơi?''

``Ơ! G-Gì đấy?''

Korone giương cặp mắt long lanh như đang xin xỏ. ``Xin phép anh cho em với Shion-chan đi chơi một buổi nha\~{}''

Kuon ngập ngừng, rồi vội vàng gật đầu: ``Ờ, ờm, được. Nhưng mà nhớ về trước ba giờ chiều đấy.''

``Vâng\~{}''

Cô khoác vai người bạn ``Shion'' rồi cùng rảo bước ra khỏi văn phòng, dáng đi hồn nhiên như không có gì xảy ra. Kuon đứng trong phòng, nhìn theo từ khung cửa sổ.

Hai bóng người đang băng qua đường. Một cô nàng năng động với tai chó lúc lắc. Một nàng tóc xám đi bên cạnh, hơi cứng nhắc, nhưng không ai nhận ra điều đó.

``\dots Cái đó, cũng chỉ là con rối của tôi thôi.''

\img{10-12i}

Cậu quay lại, thấy Shion đang đứng sát cạnh, tay chống nạnh, nháy mắt tinh nghịch.

``Em đúng là lắm trò thật đấy\dots'' cậu bật một nụ cười chát, nửa trầm trồ nửa bất lực.

``Kuon-senpai cũng đóng kịch giỏi lắm đấy chứ\~{}'' nàng Phù thủy vỗ mạnh vào lưng cậu đáp lại.

Phía trước, Korone bước đi thong thả, bỗng quay đầu lại, đôi mắt cụp xuống, nhưng môi lại cười nhẹ. ``Cuối cùng thì mình vẫn bị họ lừa được. Đúng là cái kiểu bày trò của bà\dots\ Shion-chan ạ.''

Nhưng dù sao, trò chơi khăm này, cũng đã khiến cô vui thật lòng.

%==========================================TRUYỆN PHỤ=========================================
\omk

Trong văn phòng bây giờ chỉ còn lại Kuon và Shion đang lim dim tựa vào bàn, tay chống lêm má, miệng nở một nụ cười mãn nguyện.

Cậu Tổng đứng bên cửa sổ, nhìn theo hai bóng người băng qua đường dưới ánh nắng dịu nhẹ cuối thu, nhưng rồi cậu vẫn còn đôi chút hoài nghi.

``\dots Em ấy thực sự không nhận ra luôn à?'' cậu ngoái nhìn về phía Shion, trán lấm tấm mồ hôi.

``Thì anh vừa thấy đấy. Miếng dán của em đỉnh đến mức nào nào\~{}?''

Cậu đứng khoanh tay lại, suy ngẫm lại sự việc một giây, rồi khẽ lắc đầu: ``\dots Anh không nghĩ là Koro-chan không biết đâu. Em ấy biết thừa. Chỉ là không nói thôi.''

Shion nháy mắt: ``Có thể. Nhưng chính vì thế mới vui\~{}''

Kuon chỉ thở dài lấy một hơi, bó tay với cô nàng hống hách nghịch ngợm này. Cậu ngồi phịch xuống ghế, mắt đảo qua bốn năm con rối đang ngồi xếp hàng trên sàn như những bức tượng gỗ.

Cậu lại nhìn xuống bản sao đang bất động, bản ``Shion'' mà Korone vừa trò chuyện lúc sáng.

Kế hoạch của nàng Phù thủy nghe tuy có vẻ rối rắm nhưng thật ra được thực hiện rất trơn tru. Đại ý như thế này:

\begin{itemize}
  \item Ngay khi cả hai nghe thấy tiếng bước chân của Korone ngoài hành lang, Shion đã tạo ra một bản sao đầu tiên (gọi là Bản A), trông giống hệt cô, và ngay lập tức dán miếng điều khiển lên gáy nó.
  \item Bản A sau đó cùng với Kuon được ``kéo lên trần nhà'', thực chất chỉ để đóng giả Shion thật, nhằm khiến Korone tưởng cô đã rút khỏi hiện trường. Bản sao cũng sẽ giả vờ tạo phép để khiến cho Kuon không tỏ vẻ nghi ngờ.
  \item Đồng thời, Shion thật đã tạo ra một bản sao thứ hai (bản B), cũng được gắn miếng dán, và đặt nó vào vị trí tiếp đón Korone từ cửa bước vào. Bản sao B là người mà Korone đã tương tác suốt buổi.
  \item Trong khi đó, Shion bản thật sẽ nấp ngay sau trạm máy tính, nơi mà ít người lui tới, từ đó điều khiển bản B từ  xa, quan sát mọi cử chỉ, biểu cảm, thậm chí còn sao chép lại khẩu khí của chính mình một cách hoàn hảo.
\end{itemize}

Và rồi sau đó là những gì xảy ra.

``Thế? Anh thấy kế hoạch của em vừa rồi thế nào?'' cô nàng cười tươi roi rói, ngón tay vẫn đang vung vẩy.

Cậu con trai gật gù, công nhận với bản kế hoạch công phu của cô: ``Có vài pha thót tim đấy. Nhất là lúc em ấy nhìn lên khe trần. Suýt nữa thì tưởng bị lộ rồi.''

Shion cười, mừng tượng ra khuôn mặt hoảng loạn của cậu khi nàng Khuyển ngước lên trên: ``Trêu tức cậu ta cũng vui, đúng chứ?''

Cậu gãi đầu đáp lại, tỏ vẻ hơi ngượng: ``Ừ thì\dots\ đúng thật,'' nhưng rồi giọng cậu tỏ vẻ nghiêm túc hơn: ``Nhưng mà đi lừa bạn bè kiểu này cũng chẳng hay ho gì đâu.''

Nụ cười của nàng Phù thủy khựng lại một giây. Tay cô rút lại khỏi mép bàn. ``Ồ\dots''

Tưởng rằng cô nàng sẽ lại bị cậu thuyết giảng một buổi về cách chấn chỉnh lại bản thân, nhưng không. Kuon nở một nụ cười, ẩn giấu trong đó có phần tính toán và mưu mẹo hơn thế. ``Cơ mà thỉnh thoảng bày trò một chút thì được. Như trêu bọn nai mới vào công ty, hay làm vài chiêu ảo thuật cho vui. Nhưng bạn bè thân thiết thì nên dừng ở mức đùa vui thôi.''

Shion ngẩng đầu lên, cười lớn như được cấp trên ``bật đèn xanh'': ``Đúng chứ? Mà chắc chắn là vui rồi\~{}!''

Kuon nhếch môi lên. Không phải vì đồng tình hoàn toàn với cô, mà vì cậu hiểu cái cảm giác muốn phá không khí nhàm chán ấy.

Thực ra, nếu không phải vì khả năng ``diễn'' được rèn luyện từ những năm sống với Ryujin chuyên tìm cách diễn trò để trốn uống rượu (mà mẹ cậu thường hay hiểu nhầm là say thật), và nếu không làm việc ở \hll, nơi mỗi ngày là một trận hỗn chiến nhẹ, thì Kuon đã không đủ bản lĩnh và sự sắc bén để tham gia vở kịch này.

Cậu biết rõ một điều: loài chó khôn không dễ bị lừa, và Korone là một điển hình. Chính vì thế, cậu mới càng phải diễn cho thật sâu. Và cậu cũng đã thành công.

\ct

Cả căn phòng giờ đã quay trở về bình thường, chỉ còn lại tiếng quạt chạy vi vu và tiếng lách cách của các bản sao. Kuon sán lại tới chỗ Shion, hỏi trêu: ``Này\dots\ Anh có thể mua một vài tấm miếng dán không?''

Cô nàng đáp gọn: ``Không bán.''

Cậu biết kiểu gì cô cũng sẽ như vậy, nhưng vẫn thử lại lần nữa, giọng có phần tỏ vẻ nghiêm túc hơn một chút: ``Anh sẽ trả gấp ba giá thị trường.''

``Cũng không,'' cô lập tức đáp lại.

``Vậy thế bao nhiêu?''

Không nói gì thêm, nàng Phù thủy rút ra một tập gồm mười miếng dán rồi dí vào trán rô của cậu, hóm hỉnh nói: ``Hàng tặng, không bán. Một tấm duy nhất. Dùng cho đúng người đấy\~{}''

Cậu tháo miếng dán ra, nhìn nó một lúc rồi cất vào ví, như thể giữ gì đó quý hiếm. Chẳng phải vì nó có sức mạnh ma thuật, mà vì nó là một trò đùa thông minh, và một món quà từ một người bạn lắm trò.

Phía bên sàn nhà, đám rối do Shion tạo ra đang xếp hàng như chuẩn bị tổng dượt cho buổi biểu diễn buổi chiều. Một vài cái bắt đầu nhúc nhích khi Shion ngoắc tay triệu hồi.

``Thôi, chơi thêm chút nữa, kẻo Korone về mà biết thì rách việc lắm đấy.''

``Được rồi. Nhưng nhớ đừng để rối của em đập vỡ bình cà ri của Robo-PON đó. Từ lần mà Haato-chan với Koro-chan quẳng ra ngoài\footnote{Xem lại {\flaregothic ÉPISODE 60}, quyển số Năm.} là nó mất đi độ bền rồi đấy,'' Kuon cất giọng nửa đùa nửa thật.

``Pha đó là tại chị ấy không nói rõ đấy chứ! Anh cũng bị lừa quả đó đấy thôi!''

``Không bạn. Từ cái mùi cà ri thoang thoảng bốc lên là anh thấy sinh nghi rồi. Tất cả là tại PON đấy chứ.''

``Chị ấy mà biết được là sẽ lại `Em không có PON!!' cho mà xem.''

Họ cùng cười.

Giữa một ngày cuối thu dần trôi tới hồi kết, những điều bình dị hiếm hoi vẫn có thể xảy ra, xen kẽ với những sự bất thường tới mức phi logic và không thể giải thích được. Nhưng chính những sự lổn nhổn đó là thứ khiến \hll\ luôn vui nhộn và chẳng bao giờ buồn tẻ.

Buổi sáng hôm ấy tưởng chừng sẽ kết thúc trong tiếng cười và vài trò đùa vô hại. Shion vẫn đang điều khiển một trong các bản sao, khiến nó lăn lộn trên sàn, mô phỏng một điệu nhảy thất bại đến mức khiến ai cũng phải chê cười.

``Nhìn này\~{} Con rối này kiểu như rối loạn thần kinh rồi đó!''

Cậu Tổng bật cười lắc đầu: ``Công nhận. Em nên đem bán cho bên phim kinh dị. Đảm bảo kiếm bộn tiền.''

Nhưng rồi, giữa những tiếng cười ấy, Shion chợt ngừng lại. Không còn tiếng cười khúc khích, không còn trò đùa. Cô nhìn bản sao trước mặt, mặt không biến sắc.

``\dots Ối chà. Mất điều khiển rồi.''

Kuon ngẩn mặt ra, nghiêng đầu nhìn con rối kia: ``Là sao?''

Cô nàng bình thản đáp lại: ``Em nói là mất điều khiển. Không nhận tín hiệu nữa.''

``Ý em là\dots\ miếng dán hết tác dụng à?''

``Không. Ma lực vẫn còn. Nhưng không kết nối được.''

Cô ngoe nguẩy ngón tay mình thêm một hồi lâu, nhưng con rối kia hầu như không có chút thay đổi nào về biểu cảm.

``Cái này\dots\ lạ đấy.''

Trước khi cậu kịp hỏi thêm, con rối bỗng xoay tròn dữ dội, hai tay quẫy loạn lên như đang cố thoát khỏi một điều gì đó. Và rồi\dots

Nó lao vào tấn công các bản sao khác.

Những con rối bị đâm, bị đẩy, loạng choạng va vào bàn ghế. Một cái ngã gục, một cái nắm tóc cái khác mà giật, còn một bản khác thì gào lên những âm thanh không rõ hình dạng, như tiếng hét nhưng không có giọng.

``C-Cái gì thế này!?'' Kuon ngã nhào ra sàn, cố gắng tránh khỏi những con rối đang điên cuồng tấn công lẫn nhau mà không rõ lý do. ``Sao chúng nó lại\dots?''

Nét mặt của nàng Phù thủy cũng bắt đầu căng lại, không còn giữ nét vênh váo như hồi nãy. ``Em đâu có biết! Lần đầu tiên thấy bọn nó tự tẩn nhau luôn đấy!''

``Đây là sự cố công nghệ hay ma thuật nổi loạn vậy!?'' cậu hét lên, rướn người kịp tránh một bản sao đang bị quăng về phía cậu.

Không ai kịp trả lời.

Cả đám nhân bản, sau khi đánh nhau loạn xạ, bỗng tụ lại như thể có thứ gì đó kéo hút chúng lại với nhau. Cả không gian như rung lên. Một tia sáng tím lóe lên giữa đống rối đang co cụm.

Rồi, chúng bay vọt lên trời. Như bị hút bởi một trường năng lượng vô hình, từng bản rối lao xuyên qua trần nhà từng tầng văn phòng, phá luôn cả mái nhà.

Từ dưới đất, Kuon chỉ kịp thấy những bóng người xám tro bay dần lên cao, để lại một vệt khói đen xì theo sau như đang cháy động cơ nhiên liệu.

  *{\large{\textbf{ẦM!!!}}}*

Một vụ nổ không quá lớn, nhưng đủ khiến bầu trời nhuốm một màu tím đậm, như thể ai đó đã đổ mực vào không gian. Những tia sét máu tím rạch ngang mây, uốn lượn theo hình xoắn ốc, và một cơn gió mạnh thốc qua các toà nhà khiến cửa kính rung bần bật.

Shion khẽ giật mình, như vừa nhớ ra một điều rất quan trọng: ``Chết thật! Em còn chưa kịp tiêu hủy phần lõi của con rối!''

``Ý em là sao chứ!?''

Cô nàng liến thoắng giải thích cho cậu, cố gắng không quá hoảng loạn nhưng đủ nhanh để không quá rối: ``Do ma lực bên trong chúng vẫn còn. Mà tụi nó bị mất kiểm soát nguồn cấp tín hiệu\dots\ nên tự động phản ứng theo bản năng ma thuật!''

``Bản năng ma thuật? Cụ thể là thế nào?''

Cô nhún vai, tiếp tục giải thích: ``Miếng dán chỉ là phương tiện kết nối thôi. Ma lực bên trong mới là thứ vận hành cơ thể. Khi tín hiệu bị ngắt, năng lượng đó không có chỗ thoát, thế là nó bung ra. Chưa kể, nếu môi trường xung quanh có nhiễu loạn từ trường, thì càng dễ phát nổ.''

``Nghe phức tạp như công nghệ máy tính vậy sao?'' cậu nhướng mày, chưa hiểu rõ lắm về cơ chế hoạt động của chúng. Nhưng rồi, cậu chợt nhìn thấy bản mặt của cô nàng không hề hiện rõ vẻ hoang mang nào, mà lại đang ``bình chân như vại.''

``Sao em có vẻ bình tĩnh vậy?''

``Thì\dots\ em cũng từng bị vài lần rồi. Chuyện nhỏ như con thỏ thôi.''

``\textbf{Từng bị vài lần!?}'' Kuon hét lên trong ngỡ ngàng. ``\textbf{Thế này mà nhỏ á!?}''

Shion chợt đổi giọng, sắc thái nghiêm túc hơn so với mọi khi, như thực sự không muốn đặt cậu vào thế nguy hiểm: ``Nhưng mà, nếu chuyện này xảy ra vào ngày thứ ba của chu kỳ ma lực, thì hậu quả sẽ thảm khốc hơn thế.''

Cô ngẩng mặt về phía cậu, ánh mắt sáng lên, không còn vẻ đùa giỡn thường thấy: ``Kuon-senpai, em muốn anh nhớ điều này. Trong ngày thứ ba, đừng ra đường nếu có gió xoáy màu tím. Vì rất có thể, có gì đó\dots'' giọng của cô trầm hẳn, càng khiến cho mọi thứ đang diễn ra càng hỗn loạn hơn, ``từ ngoài thế giới này sẽ tìm cách chen vào.''

``\dots Cái gì cơ?''

Cô vội vã trấn an cậu, nhưng vẫn giữ giọng cẩn trọng: ``Đây chỉ là lý thuyết thôi. Nhưng khi ma lực bị mất kiểm soát mà đồng thời xuất hiện cộng hưởng năng lượng giữa các chiều\dots\ thì cánh cửa liên vũ trụ có thể bị rạn.''

``Ý-Ý em là\dots''

Shion gật đầu. ``Va chạm vũ trụ. Không phải kiểu `đại chiến không gian' gì đâu, mà là\dots\ các luật lệ vật lý ở những thế giới khác nhau bắt đầu trộn lẫn. Em đọc được rằng \emph{một lỗ hổng sẽ xuất hiện, hút tất cả mọi thứ từ mỗi vũ trụ khác nhau và tụ hợp lại tại một điểm được gọi là `The Void (\textbf{Khoảng không})'.}''

Kuon nuốt nước bọt, cảm thấy sống lưng lạnh toát. Cậu ngước lên bầu trời đặc quánh màu đen một lần nữa, không kìm được mà nổi da gà.

Một câu bâng quơ thốt ra theo bản năng: ``\dots Không biết Koro-chan như thế nào nữa\dots''

Cô nàng vỗ vai cậu, nhẹ nhàng trấn an: ``Yên tâm đi. Cậu ấy là Thần Khuyển. Mấy cái này không làm gì được đâu.''

Quả thật, chỉ một lát sau, Korone đã quay về văn phòng, mặt mày hơi tèm lem bụi nhưng cơ thể không một vết thương nào. Cô ngó lên trời rồi quay sang hỏi với giọng nửa lo lắng nửa vui tươi: ``Tôi thấy bản sao của Shion-tan tự bay lên trời rồi nổ tung á\~{} Thế này là sao vậy?''

Nàng Phù thủy cười nhẹ, giơ tay lên đáp: ``Do tôi làm đấy.''

``\dots Anh biết ngay là Koro-chan đoán được mà.''

Korone không hỏi thêm gì. Cô chỉ gật đầu chậm rãi, rồi đi lại rót nước. Như thể chuyện đó\dots\ cũng là một ngày bình thường trong \hll\ vậy.

\dots Hoặc ít nhất, đó là trong suy nghĩ của cậu.

\il

\pov{???}

\begin{center}
  \uline{Vũ trụ 717\\Chung cư mini - Phòng Hikawa.}
\end{center}

Hầy\dots\ Lại một ngày cuối thu nắng nực nữa. Trời ạ, trời vẫn chưa thể mát hơn được sao\dots

Tôi lết từng bước chân, thở nặng nhọc sau một ngày khuôn vác. Đúng là cơ thể tôi rất khỏe, và việc nâng cả một cỗ máy một tấn không nhằm nhò gì với tôi\dots

\dots nhưng việc ngày nào cũng làm như thế đã làm cho tôi sa sút hơn hẳn cả về thể xác lẫn tinh thần.

Chưa kể đến việc kỹ năng máy tính văn phòng của tôi lại càng thui chột nữa chứ. Giá như tôi có một cái máy tính thì\dots

Chậc. Nhưng mà biết sao được. Cuộc sống của chúng tôi vốn dĩ đã được định đoạt như vậy rồi.

``\vie{Chị về rồi đây\~{}\dots}''

Cánh cửa gỗ ọp ẹp đóng lại sau lưng tôi với một tiếng cạch chán nản. Căn phòng vẫn thế, một căn phòng hẹp, cũ, và ngột ngạt. Được mỗi một cái là tinh tươm, nhưng cũng không bù nổi so với sự chật hẹp này. Mùi bụi, mùi gỗ ẩm, mùi thức ăn rẻ tiền còn chưa rửa sạch trong bồn\dots\ Tất cả chào đón tôi như một nghi thức quen thuộc.

Tôi mong chờ gì từ một căn phòng tại một chung cư mini tồi tàn cơ chứ? Của rẻ là của ôi mà, còn hơn là chẳng có chỗ để ở.

Cô em gái đang quỳ dưới đất, trán buộc một chiếc khăn, lau sàn bằng khăn ướt, tóc buộc cao nhưng vẫn rũ xuống trán. Mồ hôi nhễ nhại ứa ra, như thể Suzu đã làm việc này cả ngày hôm nay rồi.

Khi thấy tôi, em ngẩng lên cười nhẹ. ``\vie{Onee-sama về rồi ạ?}''

Tôi gật đầu, bỏ đôi giày vải đã rách gót vào góc tường, chậm rãi ngồi xuống tấm đệm đã có chút cũ kỹ. Cũng không trách được, đây là tấm đệm mà chúng tôi đã sử dụng từ hơn mười năm rồi mà.

Túi đồ ăn hôm nay chỉ có vài món giảm giá cuối ngày, thậm chí có món còn đang ở ngưỡng hết hạn sử dụng. Tôi rất muốn đi chợ và nấu một bữa ăn tử tế cho cô em gái đang tuổi ăn tuổi lớn này, nhưng sức lực của tôi còn lại sau mỗi ngày làm việc mệt nhọc không cho tôi làm việc đó. Nhìn khuôn mặt tiếc nuối vì không được thưởng thức món của tôi (nói trước là cũng rất ngon đấy), tôi không thể làm gì hơn ngoài việc luôn miệng hứa sẽ làm khi tôi rảnh, điều mà vốn dĩ cũng đã cực kỳ xa xỉ.

Dẫu vậy, em ấy vẫn vui vẻ ăn những món đồ ăn sẵn kia, không một lời kêu ca gì. Trái ngược hoàn toàn với một cô chị - là tôi đây - mệt rã rời. Cơ thể thì vậy, còn tinh thần thì rỗng không, kể từ khi bố mẹ chúng tôi lần lượt ra đi từ mười năm về trước.

``\vie{Một mình Onee-sama làm cực quá,}'' Suzu khẽ nói, giọng khàn vì bụi. ``\vie{Hay\dots\ em cũng đi làm thêm nhé? Em cũng muốn đỡ đần chị đôi chút\dots}''

Tôi nhìn đôi mắt long lanh - một màu đỏ, một màu xanh - của em ấy, khẽ lắc đầu từ chối: ``\vie{Không cần. Cứ tập trung vào học hành đi. Một mình chị gánh là được rồi.}'

Suzu nghe vậy, tỏ vẻ hơi thất vọng đôi chút. ``\vie{Nhưng mà, chị làm việc khổ cực thế kia mà!}''

``\vie{Chị không muốn em tiếp bước theo con đường cực nhọc của chị đâu,}'' tôi đáp gọn. ``\vie{Cứ tập trung học hành cho tốt, sau này kiếm được việc gì đó ổn định. Không cần phải lo cho chị đâu.}''

``\vie{Vâng ạ\~{}}'' Suzu đáp, miệng gặm một mẩu bánh mì, ngấu nghiến ngon lành.

Im lặng một lát. Chỉ còn tiếng đồng hồ cũ kỹ tích tắc vang vọng trong không gian nhỏ bé này. Tôi tự hỏi, đây cũng là lần thứ en-nờ trong suốt hơn mười năm đó, rằng liệu có khi nào\dots\ chúng tôi mới thoát khỏi tình cảnh bi đát này.

Không một người thân nào. Không một chỗ dựa. Cũng không có bạn bè nào, với tính khí nóng nảy của tôi kể từ khi chúng tôi bị dày vò bởi lũ bắt nạt trên trường, đặc biệt là lũ con trai.

Và đó cũng chính là lý do vì sao tôi cực ghét đàn ông. Không phải bởi vì tôi kỳ thị gì, chỉ là tôi hầu như không có cảm tình một chút nào. Đặc biệt là cái thằng cha chủ xưởng mà tôi đang làm kia, hắn ta chèn ép tôi tới mức quá thể luôn. Nhiều khi chỉ muốn vùng lên đấm cho hắn ta vài cái.

Cuộc đời của chúng tôi diễn ra hàng ngày chán ngắt như vậy, một cô chị đã bước sang tuổi hai mươi lăm, một cô em gái mới chỉ mười lăm tuổi.

Chúng tôi chỉ có thể cố giữ lấy nhau để sống sót - mà phải nói là sinh tồn - trong cuộc sống cực khổ tại thế gian này.

Và tôi e rằng, cả đời tôi có lẽ hai chị em chúng tôi sẽ sống như vậy, sẽ mãi nghèo rớt mùng tơi (thậm chí còn nát hơn) như thế.

\begin{center}
  \pov{-}
  \uline{Vũ trụ 256}
\end{center}

Không có gì cả.

Không ánh sáng. Không hình khối. Không chiều sâu. Không có một vật thể sống.

Chỉ là một khoảng không đen đặc, không phương hướng, không thời gian.

Một thực tại bị lột trần khỏi mọi quy luật, như thể vũ trụ này chưa từng được dựng nên từ bất kỳ công thức phép toán nào mà ta từng biết, nếu không muốn nói là tồn tại theo bất kỳ dạng thể nào.

Và rồi, từ nơi tận cùng của bóng tối, một giọng nói trầm vang lên. Nó không cần hình hài. Không cần miệng, không cần thanh quản. Chỉ cần tồn tại là quá  đủ.

``\eng{Phù\dots\ Lại đuổi thêm một người chơi khác rồi.}''

Pha trộn giữa tiếng Anh và một chút âm sắc Pháp, giọng ấy như cào nhẹ lên bề mặt hư vô. Đều đều. Dày đặc. Tựa như một vị thần vừa đẩy cánh cửa, nhìn kẻ tò mò cuối cùng rơi xuống vực.

``\eng{Họ chẳng nghe giờ nghe mình nói. Đã dặn là ở đây không có gì rồi mà\dots}''

Không có tiếng trả lời. Dĩ nhiên. Ở đây làm gì có ai mà trả lời.

Và trong im lặng, giọng nói trầm ấy lại biến mất. Không để lại dư vang. Không còn gì ngoài bóng tối.

Trò chơi, hay chính xác hơn, vũ trụ này, lại quay về điểm ban đầu.

Không có gì cả.

Và có lẽ nơi này sẽ mãi mãi chỉ là một thứ gì đó vô định, không có lấy một thứ gì.

\begin{center}
  \pov{???}
  \uline{Vũ trụ 442}
\end{center}

Lạnh.

Cái lạnh không đến từ gió trời, mà đến từ chính sự hư vô đang gặm nhấm lấy linh hồn tôi từng chút một.

Tôi đã đứng đây bao lâu rồi? Những cánh hoa anh đào đã rụng xuống rồi lại nở rộ bao nhiêu lần trên tán cây ngàn năm kia? Tôi không còn nhớ rõ nữa. Thời gian ở nơi này dường như đã chết lặng, bị giam cầm trong một vòng lặp vĩnh cửu của sương mù và những tiếng chuông gió lạc lõng.

Mọi thứ cứ lặp đi lặp lại như một cuốn phim bị hỏng. Tiếng còi tàu xa xăm, bóng tối bao trùm lấy nhà ga, và nỗi cô đơn kéo dài đến vô tận.

Tôi nhìn xuống con đường mòn dẫn lên đền Aogami. Ngày nào cũng vậy, có một cô gái\dots\ một cô gái nhỏ bé với đôi mắt u sầu, luôn lầm lũi leo lên từng bậc thang đá mòn vẹt để quỳ xuống trước điện thờ. Tôi nghe thấy tiếng thì thầm của cô ấy, cảm nhận được nỗi đau đồng điệu đang run rẩy trong từng lời cầu nguyện gửi tới Thần linh.

Cô ấy cầu xin một sự cứu rỗi. Còn tôi\dots\ tôi chỉ khao khát một sự kết thúc.

Tôi muốn được tan biến. Tôi muốn được xé toạc khỏi cái thực tại méo mó này, khỏi những sợi xích vô hình đang trói chặt tôi vào mảnh đất bị nguyền rủa. Tôi không muốn phải nhìn thấy hoàng hôn buông xuống rồi lại bình minh thức dậy trong sự tĩnh lặng đáng sợ này thêm một lần nào nữa.

Liệu có ai nghe thấy tiếng kêu cứu thầm lặng này không? Liệu có một ai đó\dots\ một thực thể nào đó đủ mạnh mẽ để bước vào đây và phá vỡ cái vòng lặp quái ác này không?

Tôi mệt mỏi lắm rồi. Linh hồn tôi đã quá rệu rã sau hàng thế kỷ bị lãng quên.

Làm ơn\dots\ hãy cho tôi được giải thoát.

\il

Năm con người, bốn vũ trụ, bốn hoàn cảnh khác nhau.

Nhưng không ai trong số họ biết rằng, ở một thời điểm không xa, năm cá thể đến từ bốn thế giới hoàn toàn khác biệt ấy - cậu Tổng quản đang chuẩn bị bước vào mùa thu thứ ba ở công ty, một cô chị đang vật lộn mưu sinh, một cô em gái đầy hy vọng, một cô bé mắt xanh lục bảo đang chờ đợi trong tuyệt vọng, và thực thể vô hình giữa hư không - sẽ cùng đặt chân đến một nơi mà không ai ngờ tới:

\begin{center}
  \uline{Vũ trụ nơi mà thực thể ``Trò chơi'' trú ngụ.}
\end{center}

\end{document}